\chapter{Analiza cerințelor și proiectarea sistemului}

Acest capitol descrie cum am transformat fundamentele teoretice din capitolul anterior într-un model concret de cerințe și decizii de proiectare pentru WhereWeFishin. Scopul nu este descrierea codului, ci justificarea structurii sistemului prin raportare la fluxurile operaționale pe care acesta trebuie să le susțină. Analiza cerințelor și proiectarea sistemului sunt etape esențiale în ingineria software, deoarece definesc baza pe care se construiește implementarea \citep{sommerville2015,pressman2014}.

Capitolul este organizat în două direcții. Prima identifică rolurile sistemului și cerințele funcționale și non-funcționale. A doua descrie principalele decizii de proiectare și modelul de domeniu prin care aceste cerințe sunt transpuse în cod.

\section{Actorii sistemului și obiectivele operaționale}

Funcționarea WhereWeFishin presupune interacțiunea dintre patru categorii de utilizatori, fiecare cu responsabilități și drepturi diferite. Această diferențiere este importantă nu doar pentru securitate, ci și pentru corectitudinea modelului de domeniu. O platformă care deservește simultan clienți finali, personal operațional și administratori nu poate folosi un model unic de interacțiune.

Cei patru actori principali sunt:

\begin{enumerate}
  \item \textbf{utilizatorul obișnuit}, care explorează locurile de pescuit pe hartă, face rezervări, își urmărește istoricul, lasă recenzii și poate încărca fișiere media pentru analiză;
  \item \textbf{angajatul}, care lucrează într-una sau mai multe locații și are acces la un set restrâns de funcții: verificarea sesiunilor și consultarea informațiilor relevante pentru locul de muncă;
  \item \textbf{managerul}, care administrează un loc de pescuit, gestionează pontoanele, angajații și informațiile descriptive, și urmărește rezervările;
  \item \textbf{administratorul sistemului}, care are o perspectivă globală, aprobă sau respinge cereri și supraveghează întreaga platformă.
\end{enumerate}

Figura~\ref{fig:use-cases} prezintă o diagramă simplificată a cazurilor de utilizare, care evidențiază legătura dintre fiecare actor și funcțiile sale principale.

\begin{figure}[ht]
\centering
\begin{tikzpicture}[
  every node/.style={font=\small},
  actor/.style={draw, rounded corners, align=center, minimum width=2.2cm, minimum height=0.8cm, fill=gray!20},
  uc/.style={draw, ellipse, align=center, minimum width=3.2cm, minimum height=0.75cm, fill=gray!5},
  arrow/.style={-Latex}
]
\node[actor] (user) at (0, 0)    {Utilizator};
\node[actor] (emp)  at (0, -4.5) {Angajat};
\node[actor] (mgr)  at (10, 0)   {Manager};
\node[actor] (adm)  at (10, -4.5){Administrator};

\node[uc] (uc1) at (5,  0.8) {Explorare hartă};
\node[uc] (uc2) at (5, -0.5) {Rezervare ponton};
\node[uc] (uc3) at (5, -1.8) {Plată online};
\node[uc] (uc4) at (5, -3.1) {Analiză media};
\node[uc] (uc5) at (5, -4.5) {Validare sesiune QR};
\node[uc] (uc6) at (5, -5.8) {Administrare locație};
\node[uc] (uc7) at (5, -7.1) {Gestionare platformă};

\draw[arrow] (user) -- (uc1);
\draw[arrow] (user) -- (uc2);
\draw[arrow] (user) -- (uc3);
\draw[arrow] (user) -- (uc4);
\draw[arrow] (emp)  -- (uc5);
\draw[arrow] (mgr)  -- (uc1);
\draw[arrow] (mgr)  -- (uc6);
\draw[arrow] (adm)  -- (uc7);
\end{tikzpicture}
\caption{Diagrama de cazuri de utilizare (simplificată)}
\label{fig:use-cases}
\end{figure}

Această structură pe roluri generează două consecințe de proiectare. Prima este necesitatea unui mecanism de autentificare și autorizare capabil să separe clar permisiunile fiecărui actor. A doua este necesitatea modelării unor fluxuri operaționale distincte, dar interconectate: o rezervare inițiată de utilizator trebuie să poată fi verificată de angajat, iar datele administrate de manager trebuie să fie vizibile în condiții controlate pentru administrator.

\section{Cerințe funcționale}

Cerințele funcționale sunt organizate pe domenii de activitate, pentru a menține o legătură clară între nevoia reală, modelul de domeniu și componentele care trebuie să o satisfacă.

Prima categorie vizează \textbf{gestionarea identității și a accesului}. Sistemul trebuie să permită înregistrarea, autentificarea și recuperarea accesului. Platforma trebuie să distingă între roluri și să restricționeze operațiile sensibile în funcție de acestea.

A doua categorie acoperă \textbf{explorarea și administrarea locurilor de pescuit}. Utilizatorul trebuie să poată vizualiza locațiile pe hartă, să consulte detalii și să compare opțiuni. Managerul și administratorul trebuie să poată crea, actualiza și controla informațiile despre locații.

A treia categorie privește \textbf{rezervările și sesiunile de pescuit}. Sistemul trebuie să permită selectarea unui interval, verificarea disponibilității și crearea unei înregistrări persistente. Calculul costului, gestionarea stării rezervării și posibilitatea anulării sunt parte din același flux.

A patra categorie se referă la \textbf{pontoane și gestionarea personalului}. Platforma trebuie să permită definirea și administrarea subzonelor rezervabile din interiorul unei locații. Managerul trebuie să poată asocia angajați locației sale.

A cincea categorie acoperă \textbf{fluxul de aplicare pentru rolul de manager}. Un utilizator poate propune spre administrare o locație și intra într-un proces controlat de aprobare sau respingere. Administratorul decide, iar decizia sa creează sau nu o locație nouă în platformă.

A șasea categorie reunește \textbf{conținutul generat de utilizatori și administrarea contextuală}: recenzii, istoricul repopulărilor cu pești și alte informații care completează experiența utilizatorului și oferă context operațional managerului.

A șaptea categorie vizează \textbf{analiza media asistată de inteligență artificială}. Sistemul trebuie să permită încărcarea de imagini și videoclipuri, trimiterea lor spre procesare și afișarea rezultatelor interpretabile. Fluxul este asincron și include validări ale fișierelor, stări explicite de procesare și prezentarea clară a rezultatelor.

Combinarea acestor cerințe arată că aplicația este simultan un marketplace, un sistem de administrare, o componentă de geolocalizare și un modul de procesare inteligentă. Tocmai această combinație justifică o arhitectură modulară.

\section{Cerințe non-funcționale și reguli de business}

Pe lângă funcționalitățile vizibile, aplicația trebuie să respecte și un set de constrângeri care garantează consistența datelor, securitatea și stabilitatea sistemului.

\textbf{Securitatea} este prima cerință non-funcțională. Platforma trebuie să protejeze identitățile, datele și operațiile sensibile. Controlul accesului este aplicat consecvent la nivel de backend, nu doar la nivel de interfață.

\textbf{Consistența temporală a rezervărilor} este critică. Sistemul trebuie să normalizeze valorile de timp, să valideze că rezervările vizează intervale viitoare și să verifice conflictele cu sesiunile existente. Starea sesiunii — \emph{Pending}, \emph{Confirmed}, \emph{Cancelled} — trebuie urmărită explicit.

\textbf{Fiabilitatea fluxurilor operaționale} presupune că procesele secundare nu blochează fluxurile critice. Dacă trimiterea unui email eșuează, rezervarea nu trebuie anulată. O analiză media cu erori trebuie marcată ca eșuată, fără a afecta restul sistemului.

\textbf{Performanța și stabilitatea răspunsului} impun ca utilizatorul să navigheze rapid și să primească feedback clar pentru operațiile de lungă durată. Procesele costisitoare sunt delegate serviciilor dedicate și executate asincron.

\textbf{Validarea datelor de intrare} este aplicată la nivel de backend. O locație trebuie să existe înainte de a fi rezervată, un ponton trebuie să aparțină locației selectate, recenziile trebuie să respecte un interval valid de rating, iar fișierele media trebuie să respecte restricții de tip și dimensiune.

\textbf{Urmărirea clară a operațiilor} este importantă pentru procese cu impact major: aprobarea unui manager, anularea unei sesiuni sau finalizarea unei analize trebuie să aibă stări și date asociate, urmăribile în timp.

\section{Proiectarea de ansamblu a sistemului}

Pe baza cerințelor formulate, WhereWeFishin este proiectată ca un sistem modular, cu componente separate pe responsabilități. Această alegere răspunde complexității funcționale și nevoii de extindere ulterioară.

Interfața web este orientată spre experiența utilizatorului și spre orchestrarea fluxurilor de interacțiune. Ea filtrează și ghidează acțiunile, dar validările esențiale și deciziile privind accesul sunt tratate în backend. Backend-ul este nucleul logic al sistemului: implementează regulile de autorizare, validările de business, controlul fluxurilor și coordonarea serviciilor externe. Persistența este tratată separat, printr-un strat dedicat și un model relațional centrat pe entitățile domeniului.

\begin{figure}[ht]
\centering
\begin{tikzpicture}[
  node distance=1.8cm and 2.2cm,
  every node/.style={font=\small},
  box/.style={draw, rounded corners, align=center, minimum width=3.4cm, minimum height=1cm, fill=gray!10},
  arrow/.style={-Latex, thick}
]
\node[box] (frontend) {Frontend Angular\\interfață web};
\node[box, right=of frontend] (backend) {Backend ASP.NET Core\\API și reguli de business};
\node[box, below=of backend] (db) {SQL Server + EF Core\\persistența datelor};
\node[box, right=of backend] (python) {Serviciu Python\\analiză imagini și video};
\node[box, above=of backend] (external) {Servicii externe\\plăți și e-mail};

\draw[arrow] (frontend) -- (backend);
\draw[arrow] (backend) -- (db);
\draw[arrow] (backend) -- (python);
\draw[arrow] (backend) -- (external);
\end{tikzpicture}
\caption{Arhitectura de ansamblu a aplicației WhereWeFishin}
\label{fig:arhitectura-wherewefishin}
\end{figure}

Microserviciul de analiză media este separat din motive practice: procesarea imaginilor și videoclipurilor implică dependențe, resurse și timpi de execuție diferiți față de operațiile obișnuite ale aplicației. Izolarea lui reduce cuplarea și limitează impactul operațiilor costisitoare asupra API-ului principal.

Serviciile externe auxiliare — procesarea plăților, notificările — sunt integrate astfel încât să poată fi înlocuite fără a afecta structura fluxurilor de business.

\section{Modelul de domeniu și relațiile dintre entități}

Modelul de domeniu descrie obiectele principale ale sistemului și relațiile dintre ele. El trebuie să acopere simultan dimensiunea geospațială, dimensiunea tranzacțională a rezervărilor și dimensiunea operațională a administrării.

\textbf{Locul de pescuit} este entitatea centrală. Reunește informații descriptive, poziționare geografică, tarif, resurse asociate și legături cu actorii care îl administrează. Este punctul de plecare pentru majoritatea fluxurilor importante.

\textbf{Pontonul} este o resursă rezervabilă distinctă, asociată unui loc de pescuit. Introducerea lui răspunde nevoii de delimitare mai precisă a spațiului: utilizatorul nu rezervă întotdeauna întreaga locație, ci adesea o zonă concretă.

\textbf{Sesiunea de pescuit} materializează fluxul de rezervare. Leagă utilizatorul de o locație, de un ponton, de un interval temporal, de un cost și de o stare operațională. Stările explicite — \emph{Pending}, \emph{Confirmed}, \emph{Cancelled} — oferă trasabilitate și susțin regulile de business.

\textbf{Utilizatorul} apare în aproape toate fluxurile. Rolul său nu este unic, ci este diferențiat prin asocierea cu responsabilități diferite. Un utilizator poate iniția rezervări, poate lăsa recenzii sau poate deveni manager printr-un proces formalizat.

Modelul include și \textbf{aplicația de manager} — care descrie procesul prin care un utilizator poate ajunge să administreze o locație — și relația \textbf{angajat--loc de pescuit}, care permite delegarea controlată a activităților operaționale.

\textbf{Recenziile}, \textbf{repopulările cu pești} și \textbf{analizele media} completează modelul. Analizele media au un ciclu de viață special, cu stări de tip \emph{Pending}, \emph{Processing}, \emph{Completed} și \emph{Failed}, care reflectă natura asincronă a procesării.

Toate aceste entități formează un model coerent, în care obiectele principale sunt conectate prin relații care reflectă activitatea reală a platformei.

\section{Fluxuri esențiale și diagrame}

Figura~\ref{fig:flux-rezervare-validare} prezintă fluxul principal al unei rezervări, de la alegerea locului de pescuit până la verificarea prin cod QR.

\begin{figure}[ht]
\centering
\begin{tikzpicture}[
  node distance=1.1cm and 0.9cm,
  every node/.style={font=\small},
  step/.style={draw, rounded corners, align=center, minimum width=2.7cm, minimum height=0.95cm, fill=gray!10},
  arrow/.style={-Latex, thick}
]
\node[step] (s1) {Alegerea\\locului de pescuit};
\node[step, right=of s1] (s2) {Selectarea\\pontonului și intervalului};
\node[step, right=of s2] (s3) {Validare\\și plată};
\node[step, right=of s3] (s4) {Crearea sesiunii\\de pescuit};
\node[step, right=of s4] (s5) {Verificare\\prin cod QR};

\draw[arrow] (s1) -- (s2);
\draw[arrow] (s2) -- (s3);
\draw[arrow] (s3) -- (s4);
\draw[arrow] (s4) -- (s5);
\end{tikzpicture}
\caption{Fluxul principal al unei rezervări și al validării sesiunii de pescuit}
\label{fig:flux-rezervare-validare}
\end{figure}

Al doilea flux important este cel al analizei media, prezentat în figura~\ref{fig:flux-analiza-media}. Față de rezervare, care este aproape sincronă, analiza media este un flux asincron, cu stări intermediare vizibile utilizatorului.

\begin{figure}[ht]
\centering
\begin{tikzpicture}[
  node distance=1.1cm and 0.7cm,
  every node/.style={font=\small},
  step/.style={draw, rounded corners, align=center, minimum width=2.8cm, minimum height=0.95cm, fill=gray!10},
  arrow/.style={-Latex, thick}
]
\node[step] (a1) {Încărcare fișier\\(imagine / video)};
\node[step, right=of a1] (a2) {Trimitere către\\serviciul Python};
\node[step, right=of a2] (a3) {Detectare și\\urmărire obiecte};
\node[step, right=of a3] (a4) {Stocare\\rezultate};
\node[step, right=of a4] (a5) {Afișare rezultat\\în interfață};

\draw[arrow] (a1) -- (a2);
\draw[arrow] (a2) -- (a3);
\draw[arrow] (a3) -- (a4);
\draw[arrow] (a4) -- (a5);
\end{tikzpicture}
\caption{Fluxul de analiză media: de la upload până la afișarea rezultatelor}
\label{fig:flux-analiza-media}
\end{figure}

Cele două fluxuri evidențiază o diferență arhitecturală importantă: rezervarea presupune un răspuns relativ rapid (sincron), în timp ce analiza media implică o procesare costisitoare, tratată asincron, cu stări intermediare.

\section{Sinteza capitolului}

Analiza realizată în acest capitol arată că WhereWeFishin este un sistem multi-rol, orientat simultan spre utilizatorul final, administrarea locurilor de pescuit și integrarea unor servicii specializate. Cerințele funcționale acoperă explorarea locațiilor, rezervarea și validarea sesiunilor, administrarea resurselor, procesul de onboarding pentru manageri și procesarea automată a fișierelor media. Cerințele non-funcționale impun constrângeri de securitate, consistență, fiabilitate și trasabilitate.

Pe baza acestor cerințe, sistemul adoptă o arhitectură modulară cu componente separate pentru interfață, logică de business, persistență și analiză media. Modelul de domeniu pune în centru locul de pescuit, sesiunea de pescuit și utilizatorul. Această structură oferă cadrul necesar pentru capitolul următor, în care accentul trece la implementarea efectivă a platformei.
